%----------------------------------------------------------------------------------------
% Tailored CV — Daten- und IT-Manager*in für Omics-Daten
% Poetsch-Gruppe, Institut für Genetik / CECAD, Universität zu Köln
% Kennziffer TUV2605-02
%----------------------------------------------------------------------------------------

\documentclass[10pt,a4paper,sans]{moderncv}

\usepackage[utf8]{inputenc}
\usepackage[T1]{fontenc}

\moderncvstyle{classic}
\moderncvcolor{blue}

\usepackage[scale=0.78]{geometry}
\usepackage{ragged2e}

%----------------------------------------------------------------------------------------

\firstname{Kundai Farai}
\familyname{Sachikonye}

\address{Auf der Lichtung 19}{82178 Puchheim, Germany}
\mobile{(+49) 1626386344}
\email{kundai.sachikonye@bitspark.com}
\homepage{github.com/fullscreen-triangle}{github.com/fullscreen-triangle}

%----------------------------------------------------------------------------------------

\begin{document}

\makecvtitle

%----------------------------------------------------------------------------------------
%	EDUCATION
%----------------------------------------------------------------------------------------

\section{Education}

\cventry{2019--2021}{PhD candidate, Computational Lipidomics}{Technische Universität München / Bayerisches Forschungsinstitut}{}{}{
  HPC mass spectrometry pipelines ($>$100\,GB mzML), multi-omics statistical modelling,
  federated learning across distributed clinical sites. Departed to pursue independent
  computational research.
}
\cventry{2018--2019}{MSc Computational Biology}{Universität Tübingen}{}{}{
  Sequence analysis, variant calling, RNA-seq quantification, pathway analysis;
  bioinformatics pipeline development.
}
\cventry{2014--2017}{MSc Pharmaceutical Biotechnology}{HAW Hamburg}{}{}{}
\cventry{2011--2014}{BSc Biochemistry and Cell Biology}{Jacobs University Bremen}{}{}{}

%----------------------------------------------------------------------------------------
%	EXPERIENCE
%----------------------------------------------------------------------------------------

\section{Experience}

\cventry{2021--present}{Independent AI \& Computational Biology Developer}{Bitspark GmbH}{}{}{
  Full-time development of production-grade scientific data pipelines, HPC analysis
  frameworks, and data management infrastructure. All systems publicly documented,
  validated, and deployed.
}

\cventry{2019--2021}{Computational Lipidomics}{TUM / Bayerisches Forschungsinstitut}{Munich}{}{
  MS-based lipidomics at HPC scale: peak detection, ion annotation (LIPID MAPS,
  LipidBlast, HMDB), multi-omics integration, federated processing across clinical
  sites. $>$100\,GB mzML dataset management.
}

\cventry{2017--2018}{Glycomics and Proteomics}{Universitätsklinikum Hamburg-Eppendorf}{}{}{
  Automated LC-MS/MS and MALDI-TOF pipeline for proteomics and glycomics;
  standardised SOP development for data acquisition and processing.
}

\cventry{2013}{Cancer Biomarker Discovery}{University Medical Center Groningen}{}{}{
  UPLC-MS/MS shotgun proteomics; data normalisation, statistical analysis.
}

%----------------------------------------------------------------------------------------
%	RELEVANT INFRASTRUCTURE AND DATA MANAGEMENT TOOLS
%----------------------------------------------------------------------------------------

\section{Infrastructure and Data Management Tools}

\cvitem{lavoisier}{
  \textbf{HPC Omics Data Pipeline.}
  Dual-pipeline MS analysis (numerical MS1/MS2 + CNN visual classification, PyTorch);
  memory-mapped I/O for datasets of arbitrary size; HDF5 structured output;
  Ray distributed execution across HPC nodes; 98.9\,\% NIST accuracy.
  Used for $>$100\,GB mzML datasets in lipidomics and metabolomics contexts.
  \href{https://github.com/fullscreen-triangle/lavoisier}{github.com/fullscreen-triangle/lavoisier}
}

\cvitem{pdsvm}{
  \textbf{Scientific Data Retrieval Infrastructure.}
  Purpose-Driven Shader VM: formalises metadata schema design as a retrieval
  problem rather than a convention problem. Each data object is indexed by
  three structured metadata coordinates $(S_k, S_t, S_e)\in[0,1]^3$;
  retrieval via Morton-ordered ternary trie runs in $O(k)$ time independent
  of collection size (373$\times$ speedup over linear scan at $N=10^5$).
  Directly applicable to HPC storage quota management and data lifecycle
  tracking for omics collections.
}

\cvitem{four-sided-triangle}{
  \textbf{Multi-Stage Scientific Data Pipeline.}
  8-stage RAG and analysis pipeline with Rust core (10--50$\times$ speedup);
  structured metadata at ingestion; SOPs implemented as Turbulance DSL protocol
  specs; FastAPI REST interface; Docker/Kubernetes deployment.
  \href{https://github.com/fullscreen-triangle/four-sided-triangle}{github.com/fullscreen-triangle/four-sided-triangle}
}

\cvitem{gospel}{
  \textbf{Genomic Variant and Multi-Omics Analysis.}
  $10^6$ variants/second VCF processing; RNA-seq integration; Bayesian network
  inference, Gaussian mixture models. Validated: 1000 Genomes Ph.\ 3, gnomAD v3.1.2,
  UK Biobank. Pathway analysis (KEGG, Reactome).
  \href{https://github.com/fullscreen-triangle/gospel}{github.com/fullscreen-triangle/gospel}
}

%----------------------------------------------------------------------------------------
%	TECHNICAL SKILLS
%----------------------------------------------------------------------------------------

\section{Technical Skills}

\cvitem{HPC / Infrastructure}{
  Ray, Dask, SLURM-compatible pipeline design; Docker, Kubernetes;
  memory-mapped I/O; HDF5, Parquet; quota management; parallel filesystem patterns
}

\cvitem{Languages}{
  Python (primary), R, Rust, Bash, SQL, TypeScript
}

\cvitem{Bioinformatics}{
  samtools, BWA, STAR, DESeq2, Nextflow (pipeline patterns); FASTQ/BAM/VCF/BED/bigWig;
  variant calling, RNA-seq quantification, ATAC-seq/ChIP-seq data types;
  pathway analysis (KEGG, Reactome, GSEA); LIPID MAPS, HMDB, NIST
}

\cvitem{ML / DL}{
  PyTorch; CNNs, LSTMs, Transformers; scikit-learn, XGBoost;
  knowledge distillation for lightweight inference; RAG pipelines
}

\cvitem{Data Management}{
  Metadata schema design; naming convention standardisation;
  SOP authoring and enforcement; data lifecycle management;
  documentation for mixed technical/non-technical audiences
}

%----------------------------------------------------------------------------------------
%	LANGUAGES
%----------------------------------------------------------------------------------------

\section{Languages}

\cvitemwithcomment{English}{Native / Fluent}{}
\cvitemwithcomment{German}{Conversational}{Living in Germany since 2011}

%----------------------------------------------------------------------------------------

\renewcommand{\listitemsymbol}{-~}

\end{document}
